\documentclass[11pt,a4paper]{article}
\usepackage[T1]{fontenc}
\usepackage{amsmath,amssymb,geometry,parskip,xcolor,hyperref}
\geometry{margin=2.3cm}
\hypersetup{colorlinks=true,linkcolor=blue!60!black,urlcolor=blue!60!black}
\title{Stage 4: can the execution HJB be made linear?}
\author{Derivation and applicability check}
\date{29 October 2026}
\begin{document}
\maketitle
% Build from this directory: pdflatex -interaction=nonstopmode -halt-on-error Stage4.tex
% The same derivation is included by ../WorkPlan.tex.
\setcounter{section}{3}
\section{Step 4: test Chen's method and derive the linear PDE}
\label{sec:stage4}

\subsection*{Answer and scope}

\textbf{Chen's paper does not establish an equivalent, closed linear PDE
for the equation in Step 2.} Its auxiliary-variable substitution gives a
linear equation only if we set aside the nonlinear relation defining that
variable. Below we apply the substitution to our equation and exhibit an
incorrect solution admitted by the resulting linear equation.

There is a usable alternative for this project's quadratic model. Add a
specified inventory diffusion, apply Cole--Hopf, and solve a linear
parabolic PDE. We prove that its logarithm gives the original value
function plus a function of time alone. Subtracting that function recovers
the original value; taking a spatial gradient recovers the original
optimal feedback directly. This additional argument is essential because
adding diffusion changes the PDE.

All calculations below assume the full inventory space $\mathbb R^M$,
constant $H\succ0$, $C\succeq0$, $\lambda\ge0$, unrestricted signed
controls, and a finite symmetric terminal penalty $\Phi\succeq0$.
These are the assumptions of Step 2. Exact liquidation is a separate
limiting problem.

\subsection*{Paper, exact location and notation map}

Use H.~J.~Kappen, \emph{A linear theory for control of non-linear
stochastic systems}, Physical Review Letters 95, 200201 (2005),
arXiv:physics/0411119v2 \cite{kappen2005}.
Open the version-specific PDF at
\href{https://arxiv.org/pdf/physics/0411119v2#page=2}{arXiv, PDF page 2}.

A downloaded copy is in the project at
\path{stage4/sources/Kappen2005.pdf}.

On PDF page 2, Eq.~(5) is the nonlinear stochastic HJB, Eq.~(6) is its
feedback, Eq.~(7) defines the logarithmic substitution, Eq.~(8) imposes
noise/control compatibility, and Eq.~(9) gives the linear operator used
in Eq.~(1) on page 1. Equations~(5)--(7) are in the left column;
Eqs.~(8)--(9) are near the top of the right column.

To prevent symbol collisions, label the paper's covariance by $N_K$
and its logarithmic scale by $a_K$ when translating its notation:
\begin{center}
\begin{tabular}{lll}
Quantity & Paper's notation & Our specialization\\\hline
Control & $\mathbf u$ & $-\mathbf v$\\
Uncontrolled drift & $\mathbf b$ & $0$\\
Control cost matrix & $R$ & $2H$\\
Noise covariance & $\nu_{ij}$, here $N_K$ & $2\nu H^{-1}$\\
Logarithmic scale & $\lambda$, here $a_K$ & $4\nu$\\
Running state cost & $V(\mathbf x,t)$ & $q(\mathbf x)=\lambda\mathbf x^\top C\mathbf x$\\
Terminal cost & $\phi(\mathbf x)$ & $G(\mathbf x)=\tfrac12\mathbf x^\top\Phi\mathbf x$
\end{tabular}
\end{center}
The matching condition becomes
\begin{equation}
 R N_K=a_K I,\qquad
 (2H)(2\nu H^{-1})=4\nu I.
 \label{eq:stage4-paper-map}
\end{equation}
Thus the paper's $\lambda$ is not our risk parameter $\lambda$.
Its operator $H$ is also distinct from our impact matrix $H$.
The diffusion addition and the proof that our deterministic feedback is
preserved are calculations for this project, not assumptions that the
paper proves for every deterministic control problem.

\subsection*{1. Identify the equation and change the time variable}

Write $A=H^{-1}$ and $q(\mathbf x)=\lambda\mathbf x^\top C\mathbf x$.
Here $q$ denotes the running inventory cost, not the inventory variable
used in the source paper for Step 2. The starting equation is
\begin{equation}
 \mathcal V_t-\tfrac14(\nabla\mathcal V)^\top A\nabla\mathcal V
 +q(\mathbf x)=0,\qquad
 \mathcal V(\mathbf x,T)=G(\mathbf x)
 :=\tfrac12\mathbf x^\top\Phi\mathbf x.
 \label{eq:stage4-original}
\end{equation}
The vector $\nabla\mathcal V$ contains the derivatives with respect to
inventory. Squaring these unknown derivatives makes this a nonlinear
first-order Hamilton--Jacobi equation. For two assets with diagonal $H$,
the nonlinear term is
\[
 \tfrac14(\nabla\mathcal V)^\top A\nabla\mathcal V
 =\frac{\mathcal V_{x_1}^2}{4\eta_1}
 +\frac{\mathcal V_{x_2}^2}{4\eta_2}.
\]
Covariance $C$ enters the known function $q$. It does not supply an
inventory second derivative in this reduced HJB.

Define time remaining $s=T-t$ and
$W_0(\mathbf x,s)=\mathcal V(\mathbf x,T-s)$. The subscript $0$ means
zero added inventory diffusion. The chain rule gives
$\partial_sW_0=-\mathcal V_t$, so
\begin{equation}
 \boxed{\partial_sW_0+
 \tfrac14(\nabla W_0)^\top A\nabla W_0=q(\mathbf x),
 \qquad W_0(\mathbf x,0)=G(\mathbf x).}
 \label{eq:stage4-inviscid}
\end{equation}
This time change is exact. It turns terminal data at $t=T$ into initial
data at $s=0$; it introduces no diffusion or approximation.
\textbf{Why take this step?} We know the cost at the execution horizon.
Using $s$ lets us evolve away from those known data with the usual
forward heat-equation time direction. The minus sign in the chain rule
is necessary to obtain that direction.

\subsection*{2. Apply the supplied paper's substitution}

\textbf{Source.} Chen \cite{chen2001}, Sec.~2, Eqs.~(1)--(7),
introduces an auxiliary unknown for a nonlinear product and collocates the
resulting linear equation at additional points. The paper leaves mutual
constraints between the approximation bases open on PDF p.~5. Its final
remarks on PDF p.~11 defer numerical experiments. Here the extra unknown
is a dependent field, not an extra spatial coordinate.

Following that substitution, introduce a scalar field $z$ by
\begin{equation}
 z(\mathbf x,s):=\tfrac14(\nabla W_0)^\top A\nabla W_0.
 \label{eq:stage4-closure}
\end{equation}
Then the equation becomes
\begin{equation}
 \boxed{\partial_sW_0+z=q(\mathbf x).}
 \label{eq:stage4-dlm}
\end{equation}
Equation~\eqref{eq:stage4-dlm} is linear in the pair of unknown fields
$(W_0,z)$: neither unknown multiplies the other. However, the original
problem requires both Eqs.~\eqref{eq:stage4-closure}
and~\eqref{eq:stage4-dlm}. Equation~\eqref{eq:stage4-closure} is nonlinear.
Keeping it preserves the original problem; dropping it admits additional
solutions. In this setting, a \emph{closed} system means that the equations
determine all unknown fields from the prescribed data.

To see the freedom directly, choose any smooth field $F$ with
$F(\mathbf x,0)=G(\mathbf x)$ and define
\[
 W_0=F,\qquad z=q-\partial_sF.
\]
This pair always solves the linear equation. There is no reason for its
$z$ to equal the squared gradient in Eq.~\eqref{eq:stage4-closure}.
Even the available initial value
$z(\mathbf x,0)=\tfrac14\mathbf x^\top\Phi A\Phi\mathbf x$
does not enforce that relation at later times.

\textbf{An explicit counterexample, derived here.}
Take one asset, $H=C=1$, $\lambda=1$, and $\Phi=0$. These are allowed
parameters. The original initial-value problem is
\[
 (W_0)_s+\tfrac14(W_0)_x^2=x^2,\qquad W_0(x,0)=0.
\]
The pair $\widetilde W=sx^2$, $\widetilde z=0$ satisfies
\[
 \widetilde W_s+\widetilde z=x^2,\qquad
 \widetilde W(x,0)=0,\qquad\widetilde z(x,0)=0.
\]
Thus it satisfies the proposed linear equation and both initial values.
But substitution into the original equation gives the residual
\begin{equation}
 \widetilde W_s+\tfrac14\widetilde W_x^2-x^2=s^2x^2,
 \label{eq:stage4-counterexample}
\end{equation}
which is nonzero when $s x\ne0$. The actual solution is
$W_0(x,s)=\tanh(s)x^2$, since
$\operatorname{sech}^2s+\tanh^2s=1$.
This counterexample concerns our application of the auxiliary-field
substitution; it is not a claim that every possible discretization in
Chen's paper returns this particular incorrect field.

\textbf{Why additional collocation points do not supply the missing
identity.} Suppose
$W_N=\sum_{k=1}^N\alpha_k\varphi_k$ and
$z_N=\sum_{k=1}^N\beta_k\chi_k$, with prescribed basis functions.
At a space--time point $\xi_j$, the linear equation gives
\[
 \sum_k\alpha_k\,\partial_s\varphi_k(\xi_j)
 +\sum_k\beta_k\chi_k(\xi_j)=q(\mathbf x_j).
\]
More sample points provide more linear rows. They do not insert
\[
 \sum_k\beta_k\chi_k(\xi_j)=
 \tfrac14\sum_{k,\ell}\alpha_k\alpha_\ell
 (\nabla\varphi_k(\xi_j))^\top A\nabla\varphi_\ell(\xi_j).
\]
The right side is quadratic in the coefficients $\alpha_k$. A square,
full-rank matrix or a unique least-squares answer therefore does not by
itself establish equivalence to our nonlinear PDE. A chosen basis could
impose extra restrictions, but those restrictions would need a separate
proof that they enforce this identity and approximate the desired solution.
\textbf{Why test this first?} A linear-looking equation is useful only if
its solutions recover the intended nonlinear problem. The counterexample
checks that requirement before we use a numerical linear solver.

\subsection*{3. Test the exponential substitution without changing the PDE}

Let $a>0$, and try $\psi_0=\exp(-W_0/a)$. For finite $W_0$ this
substitution is invertible because $\psi_0>0$. Its derivatives satisfy
\[
 (\psi_0)_s=-\frac{(W_0)_s}{a}\psi_0,
 \qquad\nabla W_0=-a\,\frac{\nabla\psi_0}{\psi_0}.
\]
Substituting Eq.~\eqref{eq:stage4-inviscid} yields
\begin{equation}
 (\psi_0)_s=
 \frac{a}{4}\frac{(\nabla\psi_0)^\top A\nabla\psi_0}{\psi_0}
 -\frac{q}{a}\psi_0.
 \label{eq:stage4-exponential-fails}
\end{equation}
The gradient-squared divided by $\psi_0$ is still nonlinear. The original
PDE has no second derivative whose transformation could cancel it.
This rules out this direct Cole--Hopf attempt on the unchanged equation;
it is not a theorem excluding every other transformation.
\textbf{Why show this calculation?} It identifies the term we still need
to cancel and explains why the next step introduces a second derivative.

\subsection*{4. Add compatible diffusion and derive a closed linear PDE}

Introduce a parameter $\nu>0$ and a new value function $W_\nu$. Define
$L_A f=\nabla\!\cdot(A\nabla f)
=\operatorname{tr}(A\nabla^2 f)$, using that $A$ is constant. Consider
\begin{equation}
 \boxed{(W_\nu)_s=\nu L_AW_\nu
 -\tfrac14(\nabla W_\nu)^\top A\nabla W_\nu+q,
 \qquad W_\nu(\mathbf x,0)=G(\mathbf x).}
 \label{eq:viscous}
\end{equation}
This is a different PDE. In a stochastic-control interpretation it
corresponds to inventory dynamics
$d\mathbf x=-\mathbf v\,dt+\sqrt{2\nu}\,H^{-1/2}d\mathbf B_t$,
where $\mathbf B_t$ is standard $M$-dimensional Brownian motion, with
expected running and terminal costs as in Step 2. Its diffusion generator
is $\nu L_A$. The noise has been added in inventory space; it is not the
price noise represented by $C$.

\textbf{Why this diffusion matrix and this exponential scale?}
The second derivative of an exponential contains a squared first
derivative. We use it to cancel the existing squared-gradient term.
First leave the exponential scale unspecified: $\psi=e^{-W_\nu/a}$,
where $a>0$. Substituting $W_\nu=-a\log\psi$ into
Eq.~\eqref{eq:viscous} gives
\begin{equation}
 \psi_s=\nu L_A\psi+
 \left(\frac a4-\nu\right)
 \frac{(\nabla\psi)^\top A\nabla\psi}{\psi}
 -\frac q a\psi.
 \label{eq:stage4-scale-choice}
\end{equation}
The nonlinear coefficient vanishes precisely when $a=4\nu$.
Both derivative terms must use the same matrix $A=H^{-1}$ for this
cancellation to hold in every inventory direction. An arbitrary isotropic
diffusion would generally fail when the impacts differ by asset.
For a diffusion $\operatorname{tr}(D\nabla^2W)$ the matching condition
is $D=aA/4$; we choose $D=\nu A$ and therefore $a=4\nu$.
Consequently choose
\begin{equation}
 \psi=\exp[-W_\nu/(4\nu)],\qquad W_\nu=-4\nu\log\psi.
 \label{eq:stage4-transform}
\end{equation}
To check the factor $4\nu$, put
$Q_\nu=(\nabla W_\nu)^\top A\nabla W_\nu$. The chain and product rules give
\begin{align}
 \psi_s&=-\frac{(W_\nu)_s}{4\nu}\psi,\nonumber\\
 \nabla\psi&=-\frac{\psi}{4\nu}\nabla W_\nu,\nonumber\\
 \nu L_A\psi&=\psi\left[-\frac14L_AW_\nu
                              +\frac{Q_\nu}{16\nu}\right].
 \label{eq:stage4-chain}
\end{align}
On the other hand, inserting Eq.~\eqref{eq:viscous} into $\psi_s$ gives
\[
 \psi_s=\psi\left[-\frac14L_AW_\nu+
                         \frac{Q_\nu}{16\nu}-\frac{q}{4\nu}\right].
\]
The first two terms are exactly $\nu L_A\psi$. Therefore
\begin{equation}
 \boxed{\psi_s=\nu\nabla\!\cdot(H^{-1}\nabla\psi)
 -\frac{\lambda}{4\nu}\mathbf x^\top C\mathbf x\,\psi
 =-\hat H\psi,\qquad
 \hat H=-\nu L_A+\frac{q}{4\nu}.}
 \label{eq:linearised}
\end{equation}
Both coefficients are known functions, so this is a closed, homogeneous
linear PDE in one unknown $\psi$. Linearity means that sums of solutions
solve the PDE with the corresponding summed initial data. In our problem
the initial data are fixed by
\begin{equation}
 \psi(\mathbf x,0)=\exp[-G(\mathbf x)/(4\nu)]
 =\exp[-\mathbf x^\top\Phi\mathbf x/(8\nu)].
 \label{eq:stage4-initial}
\end{equation}
\textbf{Why transform the data too?} The exponential is a change of
unknown throughout the initial-value problem. Retaining $G$ itself as
the initial condition for $\psi$ would solve a different problem and
would not recover the prescribed terminal cost.
In the full-space problem these positive Gaussian data, including
semidefinite $\Phi$, give the positive solution constructed below. We use
the bounded full-space evolution, not arbitrary solutions with uncontrolled
growth at infinity. Positivity permits the logarithm. Conversely, inserting
$W_\nu=-4\nu\log\psi$ into a smooth positive solution of
Eq.~\eqref{eq:linearised} recovers Eq.~\eqref{eq:viscous}.

\textbf{Source distinction.} The noise/control matching and exponential
transformation are supported by Kappen \cite{kappen2005}, Eqs.~(5)--(9).
In his notation the control matrix is $2H$, the noise covariance is
$2\nu H^{-1}$, and their product is $4\nu I$. His logarithmic scale is
therefore our $4\nu$, not our risk coefficient $\lambda$.
The calculations above specialize that principle to our equation;
they are not a Cole--Hopf result from Chen's paper.

\subsection*{5. Prove what the linear PDE recovers for our original model}

\textbf{Why is this proof necessary?} Linearizing the noisy auxiliary
problem does not yet solve the deterministic problem we started with.
We must show that the added term does not change the quantity we intend
to recover. Quadratic running and terminal costs suggest a quadratic
value; direct substitution below verifies that choice.

Because the running and terminal costs are quadratic, try
\[
 W_\nu(\mathbf x,s)=\tfrac12\mathbf x^\top P_s\mathbf x+c(s),
 \qquad P_s=P_s^\top.
\]
The subscript in $P_s$ labels time remaining; $P_s'$ denotes its
derivative. Direct differentiation gives
\[
 \nabla W_\nu=P_s\mathbf x,\quad
 L_AW_\nu=\operatorname{tr}(AP_s),\quad
 Q_\nu=\mathbf x^\top P_sAP_s\mathbf x.
\]
Match the quadratic and constant terms in Eq.~\eqref{eq:viscous}:
\begin{equation}
 P_s'=2\lambda C-\tfrac12P_sAP_s,\quad P_0=\Phi,
 \qquad c'(s)=\nu\operatorname{tr}(AP_s),\quad c(0)=0.
 \label{eq:stage4-quadratic}
\end{equation}
The equation for $P_s$ contains no $\nu$. It is exactly Step 2's Riccati
equation in the variable $s=T-t$. Thus $P_s=P(T-s)$ and
\begin{equation}
 W_0(\mathbf x,s)=\tfrac12\mathbf x^\top P_s\mathbf x,\qquad
 W_\nu=W_0+c(s),\qquad
 c(s)=\nu\int_0^s\operatorname{tr}(AP_r)\,dr.
 \label{eq:stage4-correction}
\end{equation}
One can verify $W_0$ directly in Eq.~\eqref{eq:stage4-inviscid}.
The diffusion adds a spatial constant because the Hessian of a quadratic
function is independent of inventory. This is why this particular model
preserves the feedback at every $\nu>0$; no $\nu\to0$ approximation is
needed for that statement.

The linear PDE has the explicit Gaussian solution
\begin{equation}
 \psi(\mathbf x,s)=\exp\left[-\frac{\mathbf x^\top P_s\mathbf x}{8\nu}
 -\frac14\int_0^s\operatorname{tr}(AP_r)\,dr\right].
 \label{eq:stage4-gaussian}
\end{equation}
Dividing its time derivative by $\psi$ gives
$-\mathbf x^\top P_s'\mathbf x/(8\nu)-c'(s)/(4\nu)$.
Dividing the right side of the linear PDE by $\psi$ gives
$-\operatorname{tr}(AP_s)/4+
\mathbf x^\top P_sAP_s\mathbf x/(16\nu)-q/(4\nu)$.
Equation~\eqref{eq:stage4-quadratic} makes these expressions equal.
This checks the PDE solution, including its scalar amplitude factor.

The original signed selling feedback and original value are consequently
\begin{equation}
 \boxed{\mathbf v^*(\mathbf x,t)
 =-2\nu H^{-1}\nabla\log\psi(\mathbf x,T-t)
 =\tfrac12H^{-1}P_{T-t}\mathbf x,}
 \label{eq:recovery}
\end{equation}
where $P_{T-t}$ means $P_s$ evaluated at $s=T-t$, equivalently
the matrix $P(t)$ of Step 2. In particular there is no change of selling
sign. The value itself requires the scalar correction
\begin{equation}
 \boxed{\mathcal V(\mathbf x,t)
 =-4\nu\log\psi(\mathbf x,T-t)-c(T-t).}
 \label{eq:stage4-value-recovery}
\end{equation}
The gradient removes $c$, but the value does not. A time-independent
positive rescaling of $\psi$ preserves the feedback; a time-dependent
rescaling generally changes its evolution equation.

\subsection*{6. Exact finite-dimensional check and a two-asset example}

\textbf{Why compute a reference answer?} The Gaussian solution gives an
exact full-space answer against which a future grid solver can be checked.
Comparing an independently integrated Riccati equation with a linear block
system also checks signs and factors without assuming that a PDE grid is
already accurate.

An independent linear calculation is available for this quadratic
problem. Solve the block system
\begin{equation}
 Y'=\tfrac12AZ,\qquad Z'=2\lambda CY,\qquad
 Y(0)=I,\quad Z(0)=\Phi.
 \label{eq:stage4-block}
\end{equation}
Where $Y$ is invertible,
$(ZY^{-1})'=2\lambda C-\tfrac12(ZY^{-1})A(ZY^{-1})$.
Hence $P_s=ZY^{-1}$, as in Step 2. For the stated positive quadratic
costs this construction has no finite-horizon singularity. Also
\[
 \frac{d}{ds}\log\det Y=\tfrac12\operatorname{tr}(AP_s),\qquad
 c(s)=2\nu\log\det Y(s),\qquad
 \psi=(\det Y)^{-1/2}
       \exp[-\mathbf x^\top P_s\mathbf x/(8\nu)].
\]
Here $\det Y(0)=1$ and it stays positive. This is an exact linear
\emph{ODE} system for the quadratic coefficients, not a general linear
PDE transformation supplied by Chen.

For a finite-penalty illustration, use
\[
 H=I_2,\quad C=\begin{pmatrix}1&0.4\\0.4&1\end{pmatrix},\quad
 \lambda=T=1,\quad\nu=0.2,\quad
 \Phi=\begin{pmatrix}2&0.3\\0.3&1\end{pmatrix}.
\]
This $\Phi$ is positive definite and $C\Phi\ne\Phi C$, so this example
does not assume simultaneous diagonalization of the two cost matrices.
Writing the known coefficients explicitly, the linear PDE is
\begin{equation}
 \boxed{\psi_s=0.2(\psi_{x_1x_1}+\psi_{x_2x_2})
 -1.25(x_1^2+0.8x_1x_2+x_2^2)\psi,}
 \label{eq:stage4-twoasset}
\end{equation}
with
$\psi(\mathbf x,0)=\exp[-(2x_1^2+0.6x_1x_2+x_2^2)/1.6]$ and
$\mathbf v^*=-0.4\nabla\log\psi$.
At time remaining $s=1$, the block solution gives
\[
 P_1\approx\begin{pmatrix}
 1.9643628285&0.4223113353\\0.4223113353&1.7879582164
 \end{pmatrix},\qquad c(1)\approx0.6939899962.
\]
For $\mathbf X=(3,1)^\top$, this gives
\[
 \mathbf v^*(\mathbf X,0)\approx
 (3.1576999103,1.5274461111)^\top,\qquad
 \mathcal V(\mathbf X,0)\approx11.0005458421.
\]
These numbers use the finite $\Phi$ above, not Step 3's hard terminal
constraint, so their costs need not agree with Step 3's table.

\textbf{Checks performed.} Evaluate the block solution through the
eigendecomposition of $C$ and compare it with direct fourth-order
Runge--Kutta integration of Eq.~\eqref{eq:stage4-quadratic} on $[0,1]$.
The comparison uses the 17 times $s=j/16$, $0\le j\le16$.
With 128 and 256 uniform integration steps, the maximum absolute entry
errors in $P_s$ are respectively $1.04\times10^{-10}$ and
$6.45\times10^{-12}$. With 4096 steps, the maximum errors in $P_s$
and $c$ are below $4\times10^{-15}$ and $2\times10^{-15}$,
respectively, in double precision. The observed reduction before rounding
error agrees with fourth-order integration. These are coefficient checks,
not a spatial-grid PDE experiment. The Gaussian substitution above
separately verifies the full-space linear PDE analytically.

\subsection*{7. A complete one-asset example of the successful construction}

Use the same allowed parameters as in the counterexample:
$H=C=\lambda=1$ and $\Phi=0$. This lets us compare the correct construction
with the failed auxiliary-field substitution on exactly the same problem.
There is no hard terminal liquidation constraint in this example.

Write $W_0(x,s)=b(s)x^2$. Substitution into
Eq.~\eqref{eq:stage4-inviscid} gives
\[
 b'(s)=1-b(s)^2,\qquad b(0)=0,\qquad b(s)=\tanh s.
\]
Indeed $b'=\operatorname{sech}^2s=1-\tanh^2s$. Since $P_s=2b(s)$,
Eq.~\eqref{eq:stage4-quadratic} then gives
\[
 c'(s)=2\nu\tanh s,\qquad
 c(s)=2\nu\log\cosh s.
\]
Consequently the positive linear-PDE solution is
\begin{equation}
 \boxed{\psi(x,s)=\frac{1}{\sqrt{\cosh s}}
 \exp\left[-\frac{x^2\tanh s}{4\nu}\right],\qquad
 \psi_s=\nu\psi_{xx}-\frac{x^2}{4\nu}\psi,\quad\psi(x,0)=1.}
 \label{eq:stage4-scalar-solution}
\end{equation}
To check every term, divide the derivatives by $\psi$:
\[
 \frac{\psi_s}{\psi}=-\tfrac12\tanh s
 -\frac{x^2}{4\nu}\operatorname{sech}^2s,\qquad
 \frac{\psi_{xx}}{\psi}=-\frac{\tanh s}{2\nu}
 +\frac{x^2\tanh^2s}{4\nu^2}.
\]
Thus $\nu\psi_{xx}/\psi-x^2/(4\nu)=\psi_s/\psi$ by
$\tanh^2s-1=-\operatorname{sech}^2s$. Finally,
\[
 -4\nu\log\psi-c(s)=x^2\tanh s=W_0(x,s),\qquad
 -2\nu\,\partial_x\log\psi=x\tanh s=v^*(x,T-s).
\]
\textbf{Why include the prefactor?} Without $(\cosh s)^{-1/2}$ the
spatial derivative still recovers the feedback, but the proposed $\psi$
does not satisfy the stated linear PDE. The prefactor records the scalar
correction needed for both the PDE and the recovered value.

\subsection*{Caveats and exit criterion}

The failure identified here is a missing nonlinear constraint in the
direct application of Chen's substitution. It does not establish a
universal impossibility theorem for linear representations of nonlinear
equations. The successful parabolic construction uses an added diffusion
and the project's quadratic structure. General nonquadratic costs,
inventory-dependent coefficients, sales-only controls, and imposed
boundaries require new arguments; their feedback need not be unchanged.
In particular, the full-space stochastic inventory process can cross an
inventory boundary. Feasibility of one deterministic trajectory does not
justify a sales-constrained PDE on an entire grid.

Keep $\Phi$ finite for the construction above. For exact liquidation one
may study $\Phi=\alpha I$ as $\alpha\to\infty$, but the initial Gaussian
tends to zero away from the origin and has vanishing mass. A heat-kernel
limit requires a specified normalization; it cannot be obtained by
declaring the limiting initial function to be zero almost everywhere.
The value correction and its terminal limit also require care.

\textbf{Status.} The paper's direct substitution has been tested and the
missing constraint demonstrated. The compatible-diffusion transformation,
its inverse, and exact recovery of the original quadratic feedback are
derived. The finite-penalty coefficient comparison has been checked.
Step 3's wider verification remains open.

\textbf{Exit criterion for numerical Step 4.} Solve
Eq.~\eqref{eq:stage4-twoasset} on specified finite domains with stated
boundary conditions, recover the feedback, and compare it with the same
finite-penalty Riccati solution. Refine the spatial/time grids and enlarge
the domain until the feedback error meets a declared tolerance. Report
that error separately from coefficient-integration error. This numerical
PDE validation remains to be done before downstream resource claims.

\begin{thebibliography}{9}
\bibitem{chen2001} W.~Chen (2001),
\emph{Direct linearization method for nonlinear PDE's and the related
kernel RBFs}, Sec.~2, Eqs.~(1)--(7), and Sec.~4.
\url{https://arxiv.org/pdf/math/0110005}.
\bibitem{kappen2005} H.~J.~Kappen (2005),
\emph{A linear theory for control of non-linear stochastic systems},
Phys.~Rev.~Lett.~95, 200201, PDF p.~2, Eqs.~(5)--(9).
\url{https://arxiv.org/pdf/physics/0411119v2#page=2}.
\end{thebibliography}
\end{document}
